\section{Applying comparison and rigidity to the tree passage law}
\label{tree:sec:multiscale}

The independent-walk construction supplies the local passage laws needed
by the common comparison theorem. We verify that interface with the
actual additive normalization of the field, and then give the additional
moment argument needed to initialize the tree model. Throughout this
section $\gamma=\sqrt2$ and the geometric volume base is $\lambda=2$.

\subsection{An actual-field reference with an exterior pin}

\begin{lemma}[Ordinary cone patches and exterior Gaussian pins]
\label{tree:lem:outside_pin_reference}
On every buffered ordinary patch compactly inside the cone's
normalization disk, and on each fixed affine image of such a patch,
the actual field and its passage observations can be transferred to a whole-plane Gaussian
reference whose fixed pinning circle is outside all the observation
buffers. The density depends only on the field, and is the same for
all kernels or quantum indices attached to that field marginal.
The transfer retains the actual additive constant and the reference
metric's quantum units.
\end{lemma}

\begin{proof}
Let $G_1$ be the whole-plane GFF with zero average on the unit circle.
The circle-average cone field restricted to $B_1(0)$ has the law
$G_1-\gamma\log|z|$, by \cite[Definition~4.10 and the following
discussion]{DMS}. Under the deterministic coordinate map $z\mapsto z/R$,
its representation in $B_R(0)$ is
\begin{equation}\label{tree:eq:cone_outside_pin}
 h^{\mathrm{cone},R}(z)
   =G_R(z)-\gamma\log|z|+(\gamma-Q)\log R,
 \qquad G_R(z)=G_1(z/R).
\end{equation}
Here $G_R$ has zero average on $\partial B_R(0)$. If
$U\Subset B_R(0)\setminus\{0\}$, the deterministic function added
to $G_R$ in \eqref{tree:eq:cone_outside_pin} agrees near $\overline U$
with a smooth function compactly supported away from $0$ and the
pinning circle. This function is a Cameron--Martin direction for
$G_R$. Thus the two actual field restrictions to $U$ have mutually
absolutely continuous laws.

For the shell tests, choose a fixed ordinary patch inside $B_1(0)$
and move it by a deterministic affine chart to the prescribed shell.
For example, $\phi(z)=1/2+z/128$ maps $B_{16}(0)$ compactly into
$B_1(0)\setminus\{0\}$. The Gaussian pullback is pinned on
$\phi^{-1}(\partial B_1(0))$, outside this entire buffer, and the
actual field is
\[
 G_1\circ\phi-\gamma\log|\phi|+Q\log(1/128).
\]
The smooth deterministic drift again has a Cameron--Martin extension
supported inside the pinning circle. This is an allowed outside-pinned
reference for the common comparison argument; its circle need not be
centered at the shell.

Different fixed exterior pins also give mutually absolutely continuous
local field laws. To see this directly, take $K\Subset W$ with both
pinning circles outside $\overline W$. For either pinned field, its
domain Markov decomposition on $W$ is an independent zero-boundary
GFF plus an exterior harmonic function. The pin imposes no constraint
on the zero-boundary part, whose extension vanishes on the pinning
circle. On $K$, each harmonic function admits a finite-energy cutoff
inside $W$. Conditional Cameron--Martin equivalence with the
zero-boundary restriction, followed by integration, proves mutual
absolute continuity. This argument compares actual field laws and
does not reset a random circle average.

The imaginary field is independent and has its stationary-phase law,
which is preserved by these positive affine maps. Intrinsic transport
in Proposition~\ref{tree:prop:zoom_transport} carries the full reference
experiment, its passage kernel, and the established metric in the same
chart. In particular the deterministic constants in
\eqref{tree:eq:cone_outside_pin} and in the affine pullback are retained.
On a fixed field restriction, if $Z=d\nu/d\mu$ is the density of one
reference relative to another, adjoining any of these same conditional
kernels gives, for every event $E$ and $T>0$,
\begin{equation}\label{tree:eq:reference_density_truncation}
 (\nu\otimes\mathsf K_i)(E)
 \le T(\mu\otimes\mathsf K_i)(E)
       +\E_\mu[Z\mathbf1_{\{Z>T\}}].
\end{equation}
The tail is independent of $i$. Hence uniform high-probability tests
transfer by first truncating this one density. Constant-shift
comparisons are used only on such buffered restrictions. No absolute
continuity is asserted on a domain observing a complete pinning circle.
Moment identities and uniform integrability are proved under the
original cone and contour laws; this density argument transfers
probability bounds, not moments.
\end{proof}

\subsection{The common passage-law interface}

\begin{lemma}[Tree input to the common comparison argument]
\label{tree:lem:common_interface}
Every moment and certificate extraction satisfies the path, local-kernel,
indexed-law, and joining conditions of Definition~\ref{common:admissible}
at $\gamma=\sqrt2$, $\lambda=2$, with a singleton order space.
These structural conditions give the annulus-transfer framework.
The varying-law closure condition holds for sequences with a common
interval modulus and uniformly integrable internal endpoint $k$th
powers. Proposition~\ref{tree:prop:comparison_families} supplies these
uniform bounds for the backward quantum families and each fixed finite
forward enlargement used below, together with a separate finite
comparison uniform over all fixed indices.
\end{lemma}

\begin{proof}
We check the hypotheses in their order of use.

\paragraph{Reference metric and path observations.}
Theorem~\ref{tree:thm:background_metric}, Proposition~\ref{tree:prop:background_locality},
and Proposition~\ref{tree:prop:background_confluence} give the continuous
local reference metric, its established covariance and Weyl rules,
fixed-pair uniqueness, and the strictly interior confluence input.
These are statements at $\gamma=\sqrt2$. In particular $Q>2$ leaves
the strict margin in the inward box-connection estimate used by the
common proof.
Proposition~\ref{tree:prop:certificate_limit} retains every finite
ordered marking, ranges and nonnegative subpath costs, with actual
open-margin realization and extraction of bounded-cost paths. Edge
costs vanish because the deterministic units tend to infinity.
Its second through fourth assertions give uniform local joining,
finite concatenation, and continuity of the infimal internal cost
$F_U$. The proof uses the independent Brownian contour coordinates and the
$\SLE_8$ cell topology; these conclusions are available before any
comparison or rigidity conclusion. Thus the Euclidean-circle version
of Lemma~\ref{comp:circle-chain} applies with this joining property.

\paragraph{Conditional kernels and quantum units.}
Lemma~\ref{tree:lem:finite_kernel_transfer},
Lemma~\ref{tree:lem:chart_reconstruction}, and
Proposition~\ref{tree:prop:locality} give compatible spatial kernels
from the actual real and imaginary fields. The directed local seams
supply the needed ordering information; after integrating the local
Poisson marks, no additional global order is an input. The order
space of Definition~\ref{common:admissible} can therefore be chosen
as $\{*\}$ with weight one. Disjoint real patches have product
kernels conditional on the complete fields, and independent trial
experiments use disjoint complete dependency disks. The countable
closed observation spaces are standard Borel, so the prescribed
product restrictions admit the stochastic completions used in the
common forcing proof.

Proposition~\ref{tree:prop:zoom_transport} gives the fixed-index law
\[
 \mathsf K^{(v)}(g,\widehat g)
     =(S_{v^{-p}})_\#\mathsf K(g+c_v,\widehat g),
 \qquad c_v=\gamma^{-1}\log v.
\]
It is the law of the actual transformed discrete experiment:
durations are multiplied by $v$, heights and offsets by $\sqrt v$,
and costs are read with denominator $B_nv^p$. The matched height
strips and selected-minimum tests transform with the same lattice.
Together with intrinsic affine transport, this is exactly
\eqref{comp:ann-index-definition}--\eqref{comp:ann-two-units}.
Lemma~\ref{tree:lem:outside_pin_reference} supplies each field-law
change under the unchanged local kernel, in its actual units.
Thus exterior-measurable index selection and the same-index
likelihood argument apply without a Weyl assertion for $F_U$.

\paragraph{Closure and annulus transfer.}
Under its common interval-modulus and uniform-integrability hypotheses,
Lemma~\ref{tree:lem:certificate_class_closed} retains the raw passage
array, independent reference experiments, and fixed time offsets
under further extraction, including a varying sequence of laws.
The moment bridge below verifies these bounds for the backward quantum
families and each fixed finite forward enlargement. The varying
sequences used in amplification and rigidity have indices tending to
$+\infty$ and retain only fixed offsets, so they lie in these controlled
families.
The continuum background is retained in its full joint coupling.
Its conditional law is independent of the entire raw array given
the contour and cuts; the fixed-marginal approximation in that lemma
and the coherent background construction of
Section~\ref{tree:sec:coherent-background} retain this property under changing
base units. Each fixed retained time scale diverges in lattice steps.
In particular the kernel identity is re-established from the
finite reference experiments in the new law, rather than inferred
from weak convergence of conditional independence.
The joining estimate passes through the same marked extraction.
The common Gaussian profile, annulus-transfer and local-field
modification arguments now apply. Their quantitative reference is
the stated zero-boundary Gaussian law; transfer to ordinary cone
patches uses \eqref{tree:eq:reference_density_truncation} only after
the desired probability threshold is fixed.
\end{proof}

\begin{lemma}[From indexed Gaussian charts to the whole cone]
\label{tree:lem:comparison_atlas}
If a fixed finite upper or lower comparison holds on the ordinary
Gaussian reference charts of every fixed dyadic indexed law, it holds
with the same numerical constant for the original whole cone law.
Conversely, a whole-cone comparison implies these indexed local
comparisons. Lower chart comparisons are intrinsic, as justified for
a positive constant by Lemma~\ref{comp:intrinsic-lower}; the zero lower
comparison is automatic. No positive lower comparison is assumed here.
\end{lemma}

\begin{proof}
For every fixed dyadic $v>0$, retain the
joint cone zoom of \eqref{tree:eq:background_joint_zoom}. On that
realization the passage costs and reference metric are both multiplied
by $v^{-p}$, and the spatial map is $z\mapsto r_v^{-1}z$.
Apply \eqref{tree:eq:simultaneous_background_zoom} to the raw array
and the full background before spatial projection or disintegration.
The coherent construction in
\eqref{tree:eq:bootstrap_background_kernel_equivariance}--\eqref{tree:eq:bootstrap_background_conditional}
then uses the prescribed measurable reembedding of that same background.
Thus Proposition~\ref{tree:prop:zoom_transport} identifies the resulting
passage kernel as the corresponding indexed kernel; it need not be
the original kernel. This is a joint experiment identity, not a
substitution of the field-dependent $r_v$ into versions known only
for each deterministic chart separately. Ordinary patches in the normalized unit disk
have the mutually equivalent actual-field laws of
Lemma~\ref{tree:lem:outside_pin_reference}.
Thus an almost-sure comparison for all these indexed Gaussian
restrictions gives the same numerical comparison on
$B_{r_v}(0)\setminus\{0\}$ in the original realization.
Take the countable volumes $v=2^m$, $m\ge0$. The levels
$c_v=\gamma^{-1}\log v$ tend to infinity, whereas
$r\mapsto h_r(0)+Q\log r$ is bounded above on every compact range
of radii and tends to $-\infty$ as $r\downarrow0$.
Hence $r_{2^m}\to\infty$. The preceding disks exhaust the plane.
Each fixed forward volume is allowed by the finite-window moment
bounds and the exact transport construction. Intersect the
probability-one identities for all these fixed volumes and all
ordinary rational charts before choosing a covering volume for
a particular realization.

Use a countable ordinary-chart atlas in these disks to make the local
comparisons simultaneous. A compact reference path has a finite
partition in a buffered chart cover away from zero. For a bounded-cost
certificate and a positive proposed lower constant, stop approximating
paths on successive exits from smaller members of such a cover. Each
completed exit costs at least that constant times a positive reference
separation, uniformly over the finite cover. Marked extraction therefore
bounds the number of completed exits and supplies finitely many local
pieces. This argument uses the approximating paths, without assuming
a continuous cost parametrization of every limiting certificate.
If it meets a shrinking neighborhood of zero, omit the portion from
its first to its last visit. The omitted endpoints approach zero;
Lemma~\ref{common:cone-mark-continuity} and the tree joining estimate make their joining
errors vanish. For a lower bound the omitted passage cost is
nonnegative. Thus the same comparisons hold across zero and on every
open confinement. Conversely, a comparison for the entire cone law
transports under each exact joint zoom and then under the inverse
ordinary-patch density. This proves both implications. The argument does not assert
absolute continuity under a random repinning.
\end{proof}


\subsection{The moment bridge to a finite upper comparison}

For a fixed extracted sequence put
\[
 q_i=\lim_j\frac{A_{j-i}}{A_j},\qquad b_i=2^{ip}q_i,
\]
whenever these finite limits have been extracted. For every fixed
backward index $i\geq0$, \eqref{tree:eq:internal_subadditivity} gives
$q_i\geq2^{-i}$. At any index with $q_i>0$, denote the rule in its own level-$i$ units by $G_i$;
its costs are $b_i^{-1}$ times the rule $\mathsf K^{(2^{-i})}$.
Internal endpoint norms are the actual limiting norms by
Lemma~\ref{tree:lem:moment_ui}.

We make the coarse tests precise.  After enlarging $M$ and decreasing
$m>0$, the test $\mathcal T(m,M)$ says, on a fixed buffered annular
band, that (i) its outer circle can be connected internally at cost at
most $M$, (ii) every crossing from its central hole to its outer
boundary has cost at least $m$, and (iii) the same upper and lower
statements hold for $D$.  Also include one connection across the band
of $D$-length at most $M$, and one passage connection across it of
cost at most $M$.  These last connections have endpoints on two
fixed separated circles, so the respective lower crossing bounds
apply to them. The bands are kept compactly inside the reference shell
$S=\{1/16<|z|<4\}$.
All upper assertions permit an arbitrarily small additional cost in
a fixed open margin.  These are measurable certificate tests by
Proposition~\ref{tree:prop:certificate_limit}.

\begin{proposition}[Comparison families]\label{tree:prop:comparison_families}
Consider a limit with the certificate and local-kernel properties of
Section~\ref{tree:sec:certificates}, obtained with the moment compactness
bound of Lemma~\ref{tree:lem:moment_compactness}.  Either of the following
assumptions is sufficient for exact-power upper endpoint moment
control, compactness of the normalized backward quantum family
($v\le1$) and each fixed finite forward enlargement, and a finite
deterministic upper comparison with $D$.
\begin{enumerate}
\item For some $K<\infty$, $q_i\le K2^{-ip}$ for every $i\ge0$.
\item There is an integer $i_0\ge0$ such that $0<q_i<\infty$ for every
$i\ge i_0$, and constants $0<m<M<\infty$ such that every own-unit law
$G_i$, $i\ge i_0$, satisfies $\mathcal T(m,M)$ with probability at least
$1-\varepsilon_0$.  Here $\varepsilon_0>0$ is chosen sufficiently
small, once, for the annulus-coverage threshold with favorable fraction
$1/4$ and this fixed geometry.  The finitely many scales preceding
$i_0$ need only have finite endpoint norms and satisfy the given
compactness bound.
\end{enumerate}
In the second case there is a deterministic $T_0<\infty$, depending
only on $m,M$ and the geometry, such that
\begin{equation}\label{tree:eq:b_comparability}
 T_0^{-2}\le b_i/b_j\le T_0^2\qquad(i,j\ge i_0).
\end{equation}
In either case, for a finite $K'$ and every $i\ge0$,
\begin{equation}\label{tree:eq:exact_power_endpoint}
 \|X_{2^{-i}}\|_k=q_i\le K'2^{-ip}.
\end{equation}
Each fixed finite forward enlargement has the corresponding
bounded-window controls in the retained base units. A change to its
own units is allowed once its limiting internal endpoint norm is
positive.
\end{proposition}

\begin{proof}
First consider the second case and fix a base level $i\ge i_0$.
At offset $\ell$, the physical rule in the base's own units, divided
by the reference length unit $2^{-\ell p}$, has law
\begin{equation}\label{tree:eq:relative_b_units}
 \frac{b_{i+\ell}}{b_i}G_{i+\ell}.
\end{equation}
This is the fixed-index transport of
Proposition~\ref{tree:prop:zoom_transport}; it precedes any limiting choice
of $\ell$.

All coefficients being divided by are positive by the second
hypothesis. The tests $\mathcal T(m,M)$ are precisely the two-sided
shell tests of Proposition~\ref{comp:relative-units}: outer-circle
connections supply the surrounding upper bound, and the two retained
crossings have endpoints subject to the opposite lower bound. The
kernel, affine, and joining hypotheses needed for annulus transfer
are the structural properties verified in
Lemma~\ref{tree:lem:common_interface}; varying-family compactness is
proved below after the norm bounds. Its volume
base is $\lambda=2$. Proposition~\ref{comp:relative-units}, including
the one-sided overlap argument of Lemma~\ref{comp:majority-ratios},
therefore proves \eqref{tree:eq:b_comparability}.

Equation~\eqref{tree:eq:b_comparability} yields
$q_i\le T_0^2 b_{i_0}2^{-ip}$ for $i\ge i_0$.
Increasing the constant covers the finite earlier set.  This proves
\eqref{tree:eq:exact_power_endpoint}; it is the hypothesis in the first
case.  Decompose an internal interval of arbitrary length $t\le1$
into dyadic intervals lying inside it, with at most two intervals at
each depth and all lengths at most $t$.  Stationarity, finite
concatenation, and then the endpoint modulus for the two remaining
tails give
\begin{equation}\label{tree:eq:exact_power_all_lengths}
 \|X_t\|_k\le \frac{2K'}{1-2^{-p}}t^p.
\end{equation}
The dyadic maximum bound and the mixed-monomial argument of
Lemmas~\ref{tree:lem:moment_compactness} and \ref{tree:lem:moment_ui} apply
uniformly to the backward quantum family $v\le1$, since division by
$v^p$ cancels the factor $v^p$ in
\eqref{tree:eq:exact_power_all_lengths}. On each bounded time window,
a fixed number of stationary pieces gives a common time modulus,
tight internal observations, and uniform integrability of their
endpoint $k$th powers. The certificate construction, with its countable
charts and local joining, then gives compactness of this backward
family. Each fixed finite forward enlargement has the same properties,
with constants allowed to depend on that enlargement. We do not assert
uniform tightness or uniform integrability as $v\to\infty$.

We spell out why this gives the needed uniform upper annular test.
In every limit of the compact backward family, two nearby points have
arbitrarily cheap joins in an open margin.  In a connected buffered
band, join finitely many points of a sufficiently fine Euclidean net
by a finite tree of compact paths in the band.  Subdivide its edges
into finitely many of these local joining neighborhoods.  Finite
concatenation gives a finite random bound for all outer-circle
connections.  Compactness makes these random bounds tight uniformly
over the family: otherwise extract laws with failure probability
bounded below at successively diverging bounds; in their limit a
finite set of certificates with strict cost and spatial slack gives
a finite bound, and the joining modulus transfers it back, a
contradiction.  The common reference metric likewise has, with
arbitrarily high probability, a positive crossing lower bound and a
finite band upper bound, by its continuity and local length property.
Choose the band in the ordinary cone patch of
Lemma~\ref{tree:lem:outside_pin_reference}. Its actual-field density
and \eqref{tree:eq:reference_density_truncation} transfer these bounds
to the fixed annulus reference.  Their constants can be chosen uniformly
over the backward zoom family.

Apply Lemma~\ref{comp:ann-coverage} with all sufficiently late
indices favorable and then Corollary~\ref{comp:annulus-comparison}.
This gives a finite deterministic upper comparison on the Gaussian
reference, with a constant determined only by the common shell-test
bounds and geometry.

The same numerical bound holds at every fixed dyadic volume, including
forward volumes. Indeed, for $v_0=2^m$, $m\in\Z$, the exact
indexed-kernel identity gives
\begin{equation}\label{tree:eq:forward_index_tail}
 \big(\mathsf K^{(2^m)}\big)^{(2^{-\ell})}
        =\mathsf K^{(2^{m-\ell})}.
\end{equation}
For $\ell\ge\max\{m,0\}$, the right side is a member of the backward indexed
family whose shell tests have the common bounds just proved.
Deleting the finite prefix $\ell<m$ changes only how fine the
annulus construction begins. The coverage lemma uses arbitrarily late
label windows, so this deletion changes neither their favorable
fraction nor the deterministic comparison ratio in
Corollary~\ref{comp:annulus-comparison}. Thus one finite upper factor
works for every fixed $m$. No convergence rate uniform in $m$ is
needed. A finite enlargement of a time window supplies the base
moment modulus and joining by a fixed number of stationary pieces;
it is not used to select a new upper factor depending on $m$.

Intersect the comparison events over the countable forward volumes.
Lemma~\ref{tree:lem:comparison_atlas} then gives the same finite upper
comparison on the original whole cone. This proves initial finiteness
before any argument involving optimal constants. Finite forward
enlargements also preserve the preceding endpoint and compactness
conclusions, with the corresponding fixed-window constants, and do
not affect the relative-unit tail argument beginning at $i_0$.
\end{proof}


\subsection{A deterministic metric multiplier}
\label{tree:sec:rigidity}

An upper comparison is understood in the open-port sense of
Definition~\ref{comp:port-upper}; the tree joining estimate makes it
an upper bound for $F_U$ between the exact limiting endpoints.
Define the deterministic optimal constants for the whole joint law by
\begin{equation}\label{tree:eq:rigidity_constants}
 \begin{aligned}
 C&=\inf\{q:q\text{ is an almost-sure upper comparison}\},\\
 c&=\sup\{q:q\text{ is an almost-sure lower comparison}\}.
 \end{aligned}
\end{equation}
The countable marked observations and arbitrary open losses make the
optimal constants themselves valid. Proposition~\ref{tree:prop:comparison_families}
gives $C<\infty$. For a law normalized by $A_j$, uniform integrability
makes the limiting unit internal endpoint norm equal to one.
Lemma~\ref{tree:lem:nonzero_local_cost} therefore excludes $C=0$:
zero upper comparison would give zero local costs and then a zero
internal endpoint norm by its disk-cell lifting argument. Consequently
\begin{equation}\label{tree:eq:rigidity_initial_comparison}
                 0\le c\le C<\infty,\qquad C>0.
\end{equation}
The same conclusion holds for an extraction, after a change of units,
of one of the controlled families furnished by
Proposition~\ref{tree:prop:comparison_families}, provided its limiting
internal endpoint norm is positive.

The local density statement of Lemma~\ref{comp:comparison-transport}
and the two implications of Lemma~\ref{tree:lem:comparison_atlas}
identify the optimal constants of the whole law with those of its
countable indexed Gaussian tests. The initial finite upper comparison
has already been proved, with one common constant, in
Proposition~\ref{tree:prop:comparison_families}.

\begin{theorem}[Law-level rigidity for the tree model]
\label{tree:thm:rigidity}
For every normalized tree law furnished by
Proposition~\ref{tree:prop:comparison_families}, and every extraction of
one of its controlled families after a change of units with positive
limiting internal endpoint norm, the optimal constants in
\eqref{tree:eq:rigidity_constants} satisfy $c=C>0$.
Almost surely, simultaneously on every open set $U$,
\[
                         F_U=C D_U.
\]
In particular the metric multiplier is deterministic.
\end{theorem}

\begin{proof}
We verify the remaining hypotheses of the common rigidity implication.
Lemma~\ref{tree:lem:common_interface} gives the path calculus, product
kernels, exact indexed laws and closure. Proposition~\ref{tree:prop:comparison_families}
gives compactness of the backward quantum family and each fixed
finite forward enlargement. These cover the sequences with indices
tending to $+\infty$ and fixed offsets used below. Its separate
forward-index-tail argument gives one finite upper comparison for
all fixed indices. The preceding paragraph transports its optimal numerical
constants to every fixed index and to the outside-pinned reference.
The comparison bounds pass through further marked extraction, by the
open-margin realization and local joining in
Proposition~\ref{tree:prop:certificate_limit}. Thus
Lemma~\ref{comp:ann-coverage} and Proposition~\ref{comp:amplification}
apply with $\lambda=2$.

If $c<C$, apply Proposition~\ref{comp:amplification} and retain its
high and low witnesses in the full tree observation array.
For each low witness retain its closed rectangular envelope, its two
opposite faces and endpoint positions, and its larger denominator
collar together with the actual internal reference metric there.
These are legitimate additional marks: the certificate construction
retains arbitrary finite subpaths and compact ranges, and the common
rectangle extraction selects a first crossing between opposite faces.
The range remains in the closed rectangle before the open realization
margin is added. Hence Proposition~\ref{comp:equality-tangent} applies
to this augmented extraction with its required closed envelope.
Its two witness events have separate positive probabilities; their
intersection is not assumed.

The conditional catalogue construction of
Proposition~\ref{comp:conditional-catalogue-availability} now uses the
same tree kernels. Its order-determination condition is vacuous because
the order space is a singleton. Disjoint complete dependency disks
supply the independent interior Gaussian fields and fresh passage
draws. The exact two-unit identity from
Proposition~\ref{tree:prop:zoom_transport} supplies the auxiliary law
at the prescribed shifted index. One outside pin is fixed before
all candidate supports, so every smooth catalogue modification is a
Cameron--Martin direction of that actual Gaussian reference.
The finite-energy cutoff and the density truncation in
Lemma~\ref{tree:lem:outside_pin_reference} justify all fixed reference
changes. Lemma~\ref{common:forcing-verification} supplies the catalogue
with its closed preserved rectangle, thin preservation patch, reference
connection bounds and observable outer shell crossing, including the
deterministic construction required by Lemma~\ref{comp:gate-forcing}
and the counts in Remark~\ref{comp:counting-checks}.
Product-kernel completion retains
the shortcut patch and exterior before the input is restricted to
a chosen hit.

These checks verify the hypotheses of Theorem~\ref{common:identification}.
Its proof applies the weighted Cameron--Martin argument of
Proposition~\ref{comp:rigidity-implication} and gives $c=C$, a contradiction. Therefore
\mbox{$c=C>0$} by \eqref{tree:eq:rigidity_initial_comparison}.
Lemma~\ref{comp:intrinsic-lower} upgrades the positive lower comparison
to $F_U\ge C D_U$; the upper port comparison and tree joining give
the reverse inequality. Countable exhaustion yields the simultaneous
statement on open sets. Finally, $C$ was defined from the entire
probability law in \eqref{tree:eq:rigidity_constants}, so it is a
deterministic constant before any diameter normalization is made.
\end{proof}
